% ============================================================================
%  10. Conclusiones y trabajo futuro
% ============================================================================
\section{Conclusiones y trabajo futuro}
\label{sec:conclusiones}

\subsection{Conclusiones}

\begin{itemize}
  \item El uso de CCP1 y CCP2 con Timer2 como base de tiempo única genera dos
        PWM de hardware de igual frecuencia y fase sin consumir tiempo de CPU,
        lo que deja el lazo principal libre para la navegación y la
        comunicación.
  \item La justificación matemática demostró que la especificación de
        \SI{1.2}{\kilo\hertz} no es alcanzable con \SI{20}{\mega\hertz} y
        prescaler 1:16; elegir $PR2 = 249$ da \SI{1.25}{\kilo\hertz} y una
        escala de 1000 pasos (0.1\,\%) que simplifica el firmware.
  \item Aplicar PWM directo en avance e invertido en reversa con el DRV8833
        permite controlar velocidad y sentido con solo dos líneas por motor, y
        el freno dinámico detiene la plataforma más rápido que la parada libre.
  \item Las rampas de aceleración, el freno previo a cada inversión y el
        desacoplo del bus de potencia atacan la causa principal de reinicios
        del microcontrolador: las caídas de voltaje por corriente de arranque.
  \item La arquitectura en capas y la selección de modo en compilación
        permitieron agregar el control Bluetooth sin modificar las funciones de
        control de motores.
\end{itemize}

\subsection{Competencias aplicadas}

Configuración de periféricos de microcontroladores (CCP, Timer2, USART,
interrupciones), electrónica de potencia para motores DC, diseño de firmware
embebido modular en C, comunicación serie inalámbrica, buenas prácticas de
integridad de señal y alimentación, simulación de circuitos en Proteus y
validación en hardware.

\subsection{Trabajo futuro}

\begin{itemize}
  \item Lazo cerrado de velocidad con encoders en las ruedas y control PI, para
        que el robot avance recto aunque los dos motores no sean idénticos.
  \item Seguridad por pérdida de enlace: detener el robot si no llegan
        comandos Bluetooth durante un tiempo definido.
  \item Giros suaves con velocidades distintas por rueda (giro con radio),
        además del giro sobre el eje.
  \item Medición del voltaje de batería con el ADC para alertar carga baja.
  \item Migración del montaje de protoboard a una PCB a partir del
        esquemático.
\end{itemize}
